\par\noindent\textbf{(1)} 界面无外加自由面电荷时要求 \(\Phi_1(0^+)=\Phi_2(0^-)\) 和 \(\varepsilon_1\partial_z\Phi_1(0^+)=\varepsilon_2\partial_z\Phi_2(0^-)\)。取上方 \(\Phi_1=\frac{1}{4\pi\varepsilon_1}(q/r_++q'/r_-)\)，\(r_\pm^2=\rho^2+(z\mp d)^2\)，下方 \(\Phi_2=q''/(4\pi\varepsilon_2r_+)\)。
\par\noindent\textbf{(2)} 由电势连续得 \((q+q')/\varepsilon_1=q''/\varepsilon_2\)，法向位移连续得 \(q-q'=q''\)。联立求 \(q'=q(\varepsilon_1-\varepsilon_2)/(\varepsilon_1+\varepsilon_2)\)，\(q''=2\varepsilon_2q/(\varepsilon_1+\varepsilon_2)\)，这同时满足无穷远衰减条件。
\par\noindent\textbf{(3)} 原电荷只受感应像场的力，沿 \(z\) 方向 \[F_z=qq'/(4\pi\varepsilon_1(2d)^2)=q^2(\varepsilon_1-\varepsilon_2)/[16\pi\varepsilon_1(\varepsilon_1+\varepsilon_2)d^2]\]。若下层介电常数大于上层，\(q'<0\)，电荷被吸向界面；其自能修正为 \(\Delta U=qq'/(16\pi\varepsilon_1d)\)。
\par\noindent\textbf{(4)} 当 \(\varepsilon_2\to\infty\) 时 \(q'\to-q\)，回复接地导体平面像电荷；\(\varepsilon_2=\varepsilon_1\) 时像电荷消失、力为零。镜像点电荷只是边值数学表示，不代表另一介质内有真实自由点电荷。
